\section{Discussion}
\label{sec:discussion}

\subsection{What the Results Do and Do Not Show}
\label{sec:discussion:interpretation}

The measurements in Section~\ref{sec:evaluation} show SFDB faster than the
KG baseline on every \emph{query} class measured, at every scale:
\rangelookupkg{} on LOOKUP,
\rangeglobalkg{} on GLOBAL, \rangecontextkg{} on CONTEXT,
\rangetemporalkg{} on bounded TEMPORAL
queries, and \rangetemporalunbkg{} on unbounded TEMPORAL queries. Insert
throughput does not fully share that pattern: SFDB is a clear,
repeatable win from $1\,000$ facts onward, but at the smallest scale
measured the two engines are close enough to parity that the direction
is not stable across independent runs
(Section~\ref{sec:eval:insert}). The storage-layer control
splits the query-class numbers into two different kinds of result: the LOOKUP,
GLOBAL, and CONTEXT advantages survive against KG-mem (\rangelookupkgmem{},
\rangeglobalkgmem{}, and \rangecontextkgmem{}) and are therefore properties
of context-indexed
atomic storage, while the temporal advantages do not survive at all ---
on both bounded and unbounded TEMPORAL the SFDB/KG-mem ratio straddles
parity across scale rather than settling on either side of it
(\rangetemporalkgmem{} and \rangetemporalunbkgmem{}) --- and are therefore
properties of answering queries without SQLite round trips --- something
any engine can do, and the control does. CONTEXT is the query class the
design is named for: both engines resolve it in time proportional to the
result size rather than to $N$ (neither does a full scan), so the
distinction here is not asymptotic but architectural --- SFDB reads the
answer directly from a single open set the topology already
pre-materialised at insert time (every fact registers into its own
context's open set and every ancestor's), while KG-mem still has to
enumerate the handful of distinct context buckets under the anchor and
union their fact-id sets. Whether that constant-factor difference is
large enough to matter empirically is exactly what
Table~\ref{tab:query} settles. This was not the
paper's original finding on the other three classes, on two separate
occasions, and CONTEXT was not measured at all until this revision. An earlier iteration of this evaluation reported GLOBAL queries
as SFDB's weak point --- up to three orders of magnitude slower than the KG
baseline --- because the query path enumerated topology open sets rather
than a flat fact index. That earlier finding was correct as a measurement
and wrong as an architectural conclusion: the open-set enumeration was an
implementation choice, not a consequence of context-indexing, and a flat
stalk index the engine already builds for LOOKUP queries answers GLOBAL
queries directly once the query path is changed to use it (a change of
roughly fifteen lines). A second, later iteration reported unbounded
TEMPORAL queries as a genuine architecture-level limitation rather than a
fixable implementation gap; that conclusion was also wrong, for the same
underlying reason, and Section~\ref{sec:discussion:temporal-unbounded}
describes the fix. We describe both earlier findings and their fixes here
rather than silently updating the numbers, because the fact that the
original measurements were real but the conclusions drawn from them were
not is itself informative: an implementation-level gap in a code path can
look, from the outside, exactly like a fundamental tradeoff until someone
checks whether a cheaper path can be built alongside the existing one.

The tradeoff this paper does report, and that survives this scrutiny, is
memory: SFDB uses \rangememkg{} more resident memory per fact than the
KG baseline, and \rangememkgmem{} more than the dict-indexed control
(Section~\ref{sec:eval:memory}). This is a structural
consequence of the design, not an implementation oversight --- every fact
is registered in roughly eight to ten open sets plus four auxiliary
indexes, and each registration holds a live reference into memory rather
than a compact encoded row. Reducing it would require either coarsening
the topology (fewer open sets per fact) or compressing the index
structures (Section~\ref{sec:future:compression}), both of which would
likely give back some of the latency advantage.

We do not claim, and the data do not support, that the sheaf-inspired
approach dominates a mature RDF store on every workload at every scale. Our
KG baseline is a research prototype implemented by the same team as the
sheaf engine and running against the same test harness, storage backend,
and canonicalisation pipeline. It is a controlled reference, not a
state-of-the-art. Any claim about how SFDB would fare against Jena TDB2 or
Virtuoso must wait for a comparable evaluation, which we identify as future
work.

\subsection{Fixing the Unbounded Temporal Range}
\label{sec:discussion:temporal-unbounded}

The bounded TEMPORAL results in Section~\ref{sec:evaluation} use a query
bounded on both ends, which is the case the year-bucketed temporal index
is built for; an open-ended query ("everything from year $X$ onward")
cannot be narrowed to a bounded set of buckets. An earlier version of this
evaluation measured exactly that gap: at $10^4$ facts, SFDB's TEMPORAL
latency degraded by roughly $40\times$ relative to the bounded case, to
the point of being slower than the KG baseline, and we reported it as an
architecture-level limitation --- "not a fifteen-line fix, because there
is no bounded candidate set to compute."

That claim does not survive scrutiny any better than the GLOBAL one did.
There is no bounded candidate set within the \emph{year-bucketed} index,
but nothing requires the fix to live in that index: a second, flat,
binary-searchable structure sorted by fact start/end timestamp, described
mechanically in Section~\ref{sec:complexity} and implemented in
\texttt{sfdb.sheaf.indexes.TemporalIndex}, resolves an open-ended query in
$O(\log M + n)$ rather than the $O(N)$ full-bucket fallback. This is the
exact fix Section~\ref{sec:future:global-index} had proposed as future
work before we implemented it; building it took an afternoon, not a
research programme. With it in place, unbounded TEMPORAL queries are
\rangetemporalunbkg{} faster on SFDB than the KG baseline
(Section~\ref{sec:eval:query}) instead of slower --- though the
storage-layer control shows that this residual margin, like the bounded
one, is mostly a storage-layer effect rather than a design win
(Section~\ref{sec:eval:control}). Both engines do comparable
per-candidate interval filtering once the candidate set is identified.
As with the GLOBAL case, the lesson is the
same: the original measurement was real, the "architectural" framing
around it was not.

\subsection{When the Design Helps}
\label{sec:discussion:sheaf-wins}

The context-indexed design helps when the workload has any of the
following properties. First, entity-anchored LOOKUP dominates the query mix:
each such query is answered in $O(1)$ hash lookups rather than
$O(k)$ triple joins, and the control shows this advantage belongs to the
design, not the storage layer. Second, facts have arity above two or
three: reification
overhead grows linearly with arity in the KG baseline, while the SFDB engine
stores each fact atomically. Third, queries scope naturally to sub-trees of
the context poset: the CONTEXT benchmark class is a direct test of this,
and SFDB is \rangecontextkg{} faster than KG and \rangecontextkgmem{}
faster than the dict-indexed control at every scale measured
(Section~\ref{sec:eval:query}) --- the topology's ancestor-open-set
registration turns sub-tree scoping into a single pre-materialised lookup
rather than a filter applied at query time. Fourth, the application requires
verifiable data integrity across representations: the canonical layer and
cross-engine verification harness make it cheap to run the same fact
stream through both engines and catch inconsistencies. Fifth, and
contrary to our own earlier expectation, unrestricted full-table scans
are not a case where the reified design has an advantage: whole-fact
retrieval via the flat stalk index beats per-fact reconstruction even
against the dict-indexed control. Sixth, the LOOKUP and GLOBAL advantage
is not an artefact of comparing against our own, possibly
under-optimised baselines: it holds against three real, independently
engineered production stores at every scale measured (Apache Jena TDB2,
Section~\ref{sec:eval:jena}; OpenLink Virtuoso,
Section~\ref{sec:eval:virtuoso}; Ontotext GraphDB,
Section~\ref{sec:eval:graphdb}), from a wide margin at small scale down
to a still-clear \lookupjenalarge{} (LOOKUP) and \globaljenalarge{}
(GLOBAL) against Jena, \lookupvirtuosolarge{} (LOOKUP) and
\globalvirtuosolarge{} (GLOBAL) against Virtuoso, and
\lookupgraphdblarge{} (LOOKUP) and \globalgraphdblarge{} (GLOBAL)
against GraphDB, at \jenalargescale{}
facts --- and all three stores' reported latency includes an HTTP round
trip SFDB's in-process call does not pay, which cuts against SFDB rather
than inflating its advantage.

\subsection{When the Reified Design Helps}
\label{sec:discussion:kg-wins}

The reified triple representation still wins in workloads that the SFDB
design was not built for. Temporal-range scans are now, with the control
in hand, the clearest example: once the storage layer is equalised,
SFDB shows no consistent advantage on either TEMPORAL class --- the
SFDB/KG-mem ratio straddles parity across scale on both bounded and
unbounded TEMPORAL rather than settling on either side of it
(Section~\ref{sec:eval:control}), so a workload dominated by range scans
over per-candidate filtering is better served by a reified design with a
comparable index, not worse. The real Jena TDB2 comparison sharpens this
further rather than contradicting it: Jena overtakes SFDB on
TEMPORAL from $10^4$ facts onward, at a ratio of
\temporaljenalarge{} (jena\,/\,SFDB latency,
below $1$ meaning Jena is faster by roughly seven times) at
\jenalargescale{} facts, despite paying the same HTTP overhead that works
against it on the other three classes (Section~\ref{sec:eval:jena}). The
Virtuoso and GraphDB comparisons each show the identical pattern
independently: both overtake SFDB on TEMPORAL from the same scale
onward, at ratios of \temporalvirtuosolarge{}
(Section~\ref{sec:eval:virtuoso}) and
\temporalgraphdblarge{}
(Section~\ref{sec:eval:graphdb}) --- none of
the three margins identical, but the same crossover, the same query
class, and the same scale threshold, from three independently
engineered optimisers. So this is not only a storage-layer effect a
simple dict-based control can recover, and it is not a property of one
vendor's query planner: it is a case where mature production query
optimisers genuinely outperform SFDB's Python-level candidate filtering
once the candidate sets reach tens of thousands of rows, specifically on
the one query class (TEMPORAL) where that outcome is not, per finding
ten (Section~\ref{sec:eval:lubm}), an artefact of SFDB's own
implementation leaving performance on the table. A workload
dominated by TEMPORAL queries at that scale is a real
production-engine-favoured case, not merely a KG-mem-favoured one. A
small, insert-heavy workload is a second
case: at the smallest scale measured, SFDB's insert cost is $1.22\times$
the KG baseline's, because the fixed cost of open-set registration has
not yet amortised (Section~\ref{sec:eval:insert}). Workloads that are essentially
binary---where facts genuinely are subject-predicate-object triples with
no context, no arity, and no provenance---get no benefit from the sheaf
structure and only pay the memory overhead of maintaining the extra
indexes documented in Section~\ref{sec:eval:memory}, for no corresponding
latency win. Any workload where memory is the binding constraint rather
than latency favours the reified design's more compact storage.
Finally, integration with the extensive existing RDF and SPARQL tooling
ecosystem is currently only possible on the KG side.

\subsection{Threats to Validity}
\label{sec:discussion:threats}

We identify four threats that a careful reader should weigh.

\begin{description}
    \item[Implementation bias.] Both in-house engines were written by the
          same
          authors, and both share code paths (serialisation,
          canonicalisation, planner infrastructure). We no longer have to
          speculate about how much this inflates the LOOKUP/GLOBAL/CONTEXT
          margin:
          Section~\ref{sec:eval:jena}, Section~\ref{sec:eval:virtuoso}, and
          Section~\ref{sec:eval:graphdb}
          compare against three real, independently engineered
          production stores (Jena TDB2, Virtuoso, and GraphDB), and SFDB
          remains faster on all three, on all three classes, at every scale, despite the comparison
          paying HTTP overhead that cuts against SFDB. The relative
          advantage on LOOKUP, GLOBAL, and CONTEXT against our own KG baseline is
          therefore an \emph{upper bound} that all three comparisons confirm
          does not collapse entirely once a real engine is substituted ---
          though the margins are smaller (as low as \lookupjenalarge{}
          and \globaljenalarge{} against Jena, \lookupvirtuosolarge{}
          and \globalvirtuosolarge{} against Virtuoso, and
          \lookupgraphdblarge{} and \globalgraphdblarge{} against
          GraphDB, at
          \jenalargescale{} facts, versus tens-to-hundreds-fold against
          our own baseline) and, on TEMPORAL, reverse
          outright once a real optimiser is doing the work instead of our
          KG-mem control's dict-indexed reification. That the reversal
          happens at the same scale threshold and on the same query
          class against three independently engineered optimisers, from
          three different organisations, is
          evidence the pattern is a property of mature production query
          engines generally rather than one vendor's specific
          implementation choices --- three stores are still not an
          exhaustive survey of production RDF engines, but the pattern is
          now considerably better supported than after one or two. That
          this is a three-class win rather than a two-class one is itself
          a corrected reading, not the original one: finding ten
          (Section~\ref{sec:eval:lubm}) found and fixed a genuine
          performance bug in SFDB's own CONTEXT resolution that had been
          quietly costing it the CONTEXT margin throughout this
          evaluation, including against all three production stores.

    \item[Dataset bias.] Every result up to this point uses synthetic
          datasets whose arity, context depth, and entity distribution we
          control. We no longer have to speculate about how much this
          matters: Section~\ref{sec:eval:wikidata} repeats the full
          benchmark against \wikidatanumfacts{} genuine Wikidata facts we
          did not construct, and Section~\ref{sec:eval:lubm} repeats it
          again against LUBM~\cite{guo2005lubm}, the standard benchmark
          workload for RDF stores, at four scales up to \lubmlargescale{}
          facts. Neither is a clean confirmation, though for different
          reasons. LOOKUP, GLOBAL, and CONTEXT reproduce the synthetic
          finding on both real datasets; TEMPORAL does not, and lands in a
          different place on each: it beats both baselines on Wikidata
          (the opposite of the synthetic-benchmark control result), and
          sits at roughly parity with both on LUBM's two largest scales
          (neither the clean win of Wikidata nor the clear loss against
          KG-mem on synthetic data). Real knowledge graphs have long
          tails, correlated structure, and adversarial cases that a
          controllable generator does not model by construction, and
          TEMPORAL landing in three different places across three
          datasets is direct evidence that one of the five
          margins this paper reports is genuinely sensitive to exactly
          that, in a way LOOKUP, GLOBAL, and CONTEXT are not. Building the
          LUBM comparison also surfaced finding ten (the bug inventory
          above): a genuine performance bug in SFDB's CONTEXT resolution
          that LUBM's topology shape (many large, mutually unrelated
          top-level contexts) exposed and neither the synthetic generator
          nor Wikidata's data shape had. Two organic-or-standard datasets
          is still not a survey; DBpedia and BSBM/WatDiv-as-workloads
          remain untried.

    \item[Query bias.] Our query set was designed with both engines in
          mind: one anchor-based class (LOOKUP), one no-anchor class
          (GLOBAL), one sub-tree class (CONTEXT), and one range-based class
          measured under two parameter
          settings (bounded and unbounded TEMPORAL). A different query mix
          would shift the aggregate numbers; Section~\ref{sec:discussion:temporal-unbounded}
          shows that varying a single parameter of the TEMPORAL class
          (whether the range is bounded) used to be enough to flip which
          engine is faster, before we added the index that closes that gap
          against the SQLite baseline. We do not claim the five classes
          measured here are exhaustive. Unlike an earlier revision of this
          paper, we also do not claim that the KG baseline's scan-and-filter
          approach never wins: against the KG-mem storage-layer control,
          bounded TEMPORAL is exactly such a case, with the SFDB/KG-mem
          ratio straddling parity across scale rather than settling on
          SFDB's side of it (Section~\ref{sec:eval:control}).

    \item[Scale.] All measurements are at $10^5$ facts or below. At the
          largest scale each fact expands to roughly $15$ reified triples
          in the KG representation, so the KG engine is handling on the
          order of $1.5$ million triples --- past the entry point of
          standard RDF benchmarks, but far below the $10^8$-triple scales
          at which production stores are routinely evaluated. Our design
          choices are consistent with the complexity analysis in
          Section~\ref{sec:complexity}, but empirically our claims are
          bounded at $10^5$ facts on a single node.
\end{description}

We also record, in the interest of showing our work rather than only its
conclusion, that ten bugs, performance gaps, and measurement findings
were surfaced
during this evaluation, by eight different mechanisms, only one of which
was the evaluation methodology catching itself automatically: the
automated cross-engine harness (once); manual code inspection noticing
an existing cheaper structure (three times); manual inspection triggered
by building an unrelated external comparison that demanded exact
semantics (once); a stress test failing outright at a larger scale than
previously exercised (once); a targeted correctness test authored to
validate an unrelated optimisation, which surfaced an independent bug as
a side effect (once); repeating a seeded, supposedly deterministic
benchmark in an independent process while building a second external
comparison, which surfaced a measurement-methodology gap rather than a
code bug (once); building a third external comparison, which
surfaced a cross-engine SPARQL semantics divergence in the benchmark's
own shared query text rather than a code bug in any engine (once); and
extending a standard benchmark workload with a topology shape (many
large, mutually unrelated top-level contexts) none of the project's own
datasets had exercised, which surfaced a genuine performance bug that
had been silently inflating a majority of this paper's own CONTEXT
measurements throughout the entire evaluation, invisible to the
correctness-checking harness because it never affected results, only
how much work produced them (once):

\begin{enumerate}
    \item The sheaf engine's LOOKUP path silently used neighbourhood
          (any-mention) match semantics instead of strict subject match.
          This one genuinely was caught by the cross-engine verification
          harness described in Section~\ref{sec:correctness}, after we
          tightened it from a positional, stringified comparison to an
          order-independent, type-preserving one: the two engines'
          canonical result sets stopped agreeing, which is exactly what
          the harness is for.
    \item The GLOBAL open-set enumeration inefficiency that produced an
          earlier draft's headline negative result (Section~\ref{sec:discussion:interpretation})
          was \emph{not} caught by the harness, because both engines
          already agreed on which facts a GLOBAL query returned before the
          fix --- they disagreed only on how long it took. It was found by
          inspecting the query path and noticing a cheaper index was
          already being maintained elsewhere for LOOKUP.
    \item A linear scan standing in for a dictionary lookup in the KG
          engine's identifier decoder was found the same way, by code
          inspection rather than by any automated check, and had been
          inflating the KG engine's own reconstruction cost on every
          query class that touches reified triples.
    \item The unbounded-TEMPORAL degradation
          (Section~\ref{sec:discussion:temporal-unbounded}) was found and
          fixed the same way as the GLOBAL case: code inspection, not the
          harness, since it changed only how long a correct answer took to
          compute.
    \item The benchmark's ``unbounded'' TEMPORAL query was, for one
          revision, not actually unbounded: both engines share a
          bare-year convenience shorthand under which a start of
          ``\texttt{2023}'' with no end bound is silently re-bounded to
          the single year $[2023, 2024)$, so the suite measured a
          re-bounded query while describing it as open-ended. The
          harness could not catch this --- both engines implemented the
          same unintended semantics and agreed perfectly --- and it was
          found by reading the shared query-parsing code while building
          the external rdflib comparison
          (Section~\ref{sec:eval:rdflib}), which needed the exact
          semantics spelled out in SPARQL. The fix issues a full ISO
          timestamp that the shorthand does not touch, and a regression
          test now pins the suite's unbounded query to a genuinely open
          end.
    \item The sheaf engine's restriction-graph builder compared every
          pair of open sets to determine subset relations, an
          $O(|\text{open sets}|^2)$ construction that was invisible at
          the $10^2$--$10^4$-fact scales this evaluation used before this
          revision --- tens of thousands of open sets at $10^5$ facts
          made it run for over half an hour without completing. This was
          found by neither the harness nor code inspection but by scaling
          up: extending the evaluation to $10^5$ facts to address the
          threat named in Section~\ref{sec:discussion:threats} (Scale)
          simply hung, which a fixed evaluation scope would never have
          surfaced. The fix replaces the pairwise sweep with a
          membership-driven construction (intersecting the topology's
          point-to-open-set map) that is linear in total open-set
          memberships; a targeted test recomputes the old pairwise result
          as an oracle and asserts the two constructions agree.
    \item Writing that test surfaced a second, unrelated bug in the same
          code path: the restriction graph is only ever grown by
          \texttt{add\_edge} and was never cleared before a rebuild, so
          an edge recorded before an insert broke the subset relationship
          it depended on would survive in the graph forever. This does
          not affect any number reported in Section~\ref{sec:evaluation}
          --- the benchmark inserts a fact stream once and only then
          issues queries against a fixed topology, so no run ever
          triggers a rebuild after the topology has changed since the
          previous one --- but it is a genuine correctness gap in a code
          path SEMI-LOCAL query routing depends on for any workload that
          interleaves inserts and queries, which this evaluation's own
          insert-once-then-query-many pattern cannot exercise. We fixed
          it (a \texttt{RestrictionGraph.clear()} called before every
          rebuild) and added a direct regression test rather than treat
          "the benchmark doesn't hit it" as grounds to leave it.
    \item Building the Apache Jena TDB2 comparison (Section~\ref{sec:eval:jena})
          required repeating the $10^5$-fact CONTEXT benchmark in a
          separate process from the one that produced
          Section~\ref{sec:eval:query}'s numbers, to run it alongside the
          Jena queries under identical conditions. The two measurements
          disagreed by $\sim\!45\%$ ($712$\,ms versus $1\,044$\,ms) despite
          both being ten-sample means of a seeded, deterministic benchmark
          with sub-$2\%$ within-run confidence intervals. This is not a
          code bug --- both numbers are real measurements of the same
          code --- but a gap in what the within-run statistics can see:
          they capture noise \emph{within} one process's ten samples, not
          variation \emph{across} independent process launches. We
          reproduced the higher figure with no other process running at
          all, ruling out contention with the concurrent Jena/Fuseki
          server as the explanation. A third independent process launch,
          needed to apply the query-portability fix in the next finding,
          produced a third value ($989$\,ms) in the same span,
          reinforcing rather than undermining the finding.
    \item Building the Virtuoso comparison (Section~\ref{sec:eval:virtuoso})
          surfaced a divergence not in SFDB's code but in how the three
          external engines interpret the shared SPARQL query text
          itself. The TEMPORAL and unbounded TEMPORAL query text
          compares an RDF literal against another literal with
          \texttt{<} or \texttt{>}; against Virtuoso this silently
          returned every row unfiltered rather than a spec-consistent
          comparison, while the identical query text against Jena and
          rdflib correctly excluded the rows it should have and agreed
          with SFDB. We traced this to how Virtuoso resolves relational
          comparisons over untyped RDF literals and fixed it once, in
          the shared query text used by all three external comparisons,
          by wrapping the compared variable in \texttt{STR(...)} to
          force an unambiguous string comparison; all three external
          comparisons were re-verified to still agree with SFDB after
          the change. Unlike most of the findings in this list, this one
          is
          not a bug in code this project controls: it is a reminder that
          identical query text is necessary but not sufficient for a
          fair cross-engine comparison when the engines themselves
          disagree about SPARQL's own semantics at the margins.
    \item Building the LUBM standard-benchmark-workload comparison
          (Section~\ref{sec:eval:lubm}) surfaced the most consequential
          finding in this list: a genuine performance bug in SFDB's
          CONTEXT query resolution, present for the entire evaluation and
          invisible to every prior comparison. LUBM's largest scale
          generates $265$ mutually unrelated top-level contexts (one per
          university); a CONTEXT query anchored on one small, fixed
          department went from $0.68$\,ms to $350$\,ms as unrelated
          universities were added elsewhere, a $515\times$ slowdown with
          the query's own result size held constant at $68$ facts the
          entire time. The cause
          (\texttt{src/sfdb/sheaf/optimizer.py},
          \texttt{\_classify\_semilocal}): CONTEXT resolution found every
          fact whose exact context matched the query and unioned in
          every open set each belonged to, including the open set named
          for \texttt{context:world} --- which every fact in the corpus
          belongs to whenever the query context is not the root ---
          instead of taking the direct $O(1)$ lookup on the open set
          named \texttt{context:<c>}, which the architecture already
          materialises at insert time and which
          Section~\ref{sec:complexity}'s complexity analysis already
          assumed was what CONTEXT resolution did. Final results were
          always correct (a downstream filter narrows the over-broad
          candidate set back down before it reaches the caller), which
          is exactly why cross-engine verification never caught it across
          nine prior findings' worth of scrutiny, and why the balanced
          synthetic generator never exposed it either: that generator
          grows a query's own target subtree in lockstep with total
          corpus size, so the wasted work and the legitimate work scale
          together and look like the same line. Only a topology shape
          with many large, mutually unrelated top-level siblings ---
          which none of this paper's other datasets have --- decouples
          them. The fix targets \texttt{context:<c>} directly;
          \texttt{tests/sheaf/test\_context\_query\_scope.py} pins it,
          confirmed to fail against the pre-fix code. Because CONTEXT
          queries appear in every comparison this paper reports, this
          finding required re-running and re-verifying every benchmark in
          Section~\ref{sec:evaluation} rather than a local patch; two
          earlier claims in this paper (the eighth finding above, and a
          CONTEXT-loses-to-KG-mem reading in
          Section~\ref{sec:eval:wikidata}) are retracted as a direct
          consequence, and the corrected numbers show SFDB winning
          CONTEXT robustly everywhere it is measured rather than losing
          it from $10^4$ facts onward against three production stores.
          We report the retractions alongside the fix rather than quietly
          substitute corrected numbers into the earlier sections and
          leave the narrative as though CONTEXT had always been this
          fast. The fix was subsequently re-verified under an even
          harder test than any single comparison in this paper already
          ran: LUBM data against all three external stores at once
          (Section~\ref{sec:eval:lubm-external}), varying the dataset and
          the comparison engine simultaneously. The corrected CONTEXT
          advantage held there too, at every scale, against every store.
\end{enumerate}

All ten findings are reflected in the numbers reported in
Section~\ref{sec:evaluation}, Section~\ref{sec:eval:jena},
Section~\ref{sec:eval:virtuoso}, Section~\ref{sec:eval:graphdb}, and
Section~\ref{sec:eval:lubm}. That only the first was actually caught by
an automated check is itself a threat to validity worth naming plainly:
the verification harness is real and it does its job, but it checks
result equivalence, not performance or intent, and it cannot catch
an inefficiency both engines agree on, a specification bug
both engines implement identically, or a correctness gap in a code path
the benchmark's own access pattern never exercises --- the tenth finding
is the sharpest instance of exactly that last case: a bug that silently
inflated a majority of the paper's own CONTEXT numbers throughout the
whole evaluation, past nine prior rounds of scrutiny, because it never
touched a single returned result. The sixth finding adds a
further point: a fixed evaluation scope can hide a scalability cliff that
only appears once you actually push past it, which is itself an argument
for the $10^5$-fact scale this revision adds over the $10^4$ ceiling of
the one before it. The eighth, as originally stated, claimed something
different: that even a seeded, deterministic benchmark is not
automatically reproducible \emph{across} process launches to the
precision its own within-run confidence intervals suggest. We no longer
stand behind that specific claim --- Section~\ref{sec:eval:jena} retracts
it, because the numbers it rested on turned out to be measurements of
the bug the tenth finding fixes, not of genuine process-level variance
--- and we are explicit about that rather than let a wrong methodological
lesson stand uncorrected alongside the finding that undermines it. The
ninth
adds a point that does still stand: "identical query text across
engines" is a methodological
commitment worth stating plainly (Section~\ref{sec:eval:jena}), but it
is not by itself a guarantee of a fair comparison, since the engines
being compared can silently disagree about what that text means. The
tenth adds the sharpest point of all: a bug that never affects a single
returned result can still be a bug worth an entire evaluation's
recomputation, and "the harness verifies correctness" is not the same
claim as "the harness verifies the paper's numbers are what they should
be."

\subsection{Negative Results}
\label{sec:discussion:negative}

Four negative results deserve to be stated plainly, because they
discipline the paper's overall claim. (An earlier version of this
section reported a different one in place of the second below --- TEMPORAL
degradation under an
open-ended range --- which Section~\ref{sec:discussion:temporal-unbounded}
describes fixing; we note that history rather than silently rewrite it.)

The first is the memory cost
quantified in Section~\ref{sec:eval:memory}: SFDB uses \rangememkg{}
more resident memory per fact than the KG baseline and \rangememkgmem{}
more than the dict-indexed control. This is the tradeoff
that funds every latency advantage reported in this paper, and it is
intrinsic to the design rather than an implementation artefact: redundant
per-open-set registration is what makes the stalk and temporal indexes
cheap to query.

The second is new in this revision and comes directly from the
storage-layer control: the TEMPORAL advantages over the SQLite baseline
do not survive against KG-mem on either class --- on both bounded and
unbounded TEMPORAL the SFDB/KG-mem ratio straddles parity across scale
rather than settling on either side of it
(\rangetemporalkgmem{} bounded, \rangetemporalunbkgmem{} unbounded).
Context-indexed storage buys nothing measurable for range scans whose
cost is dominated by per-candidate filtering; the earlier draft's
temporal ``wins'' were storage-layer effects. We consider surfacing this
worth more than the wins it retracts: it is precisely the kind of
finding the control exists to produce.

The third, also new in this revision, is that insert throughput does not
favour SFDB unconditionally: at the smallest scale measured ($100$
facts) the two engines are close enough to parity
(\insertkgsmallestscale{} in the reported run) that the direction is
not stable across independent process launches --- repeated launches of
the identical measurement put the ratio on both sides of parity --- and
SFDB only becomes a clear, repeatable win from $1\,000$ facts onward
(Section~\ref{sec:eval:insert}). The fixed per-insert cost of registering
a fact into eight to ten open sets does not amortise until the fact
stream is large enough, which means the design's insert advantage is a
statement about scale, not an unconditional property of the engine.

The fourth is the vacuousness of the sheaf condition in our setting
(Section~\ref{sec:sheaf-remark}): the system is, strictly, a presheaf
database with a defensive consistency check, not a sheaf database in any
technical sense that would constitute a contribution on its own. This does
not diminish the practical value of the design: what makes the design work
is the context poset, the stalk index, the canonical intermediate
representation, and the disciplined restriction vocabulary that comes with
them. It does mean that any framing of the paper
as a sheaf-theoretic breakthrough would be dishonest, and we do not adopt
that framing.
